\documentclass[11pt]{article}
\usepackage[a4paper,margin=25mm]{geometry}
\usepackage{amsmath,amssymb,amsthm}
\usepackage[colorlinks=true,urlcolor=blue]{hyperref}
\newtheorem{lemma}{Lemma}
\newtheorem{proposition}{Proposition}
\newtheorem{theorem}{Theorem}
\newcommand{\OPT}{\operatorname{OPT}}
\title{A tight guarantee for iterative rounding\\in the two-size Gasoline problem\\\large Mathematical proof and computational checks}
\author{MathIsEvenEasier}
\date{4 October 2026}
\begin{document}
\maketitle
\noindent\textbf{Status.} This document gives a mathematical proof for positive
integer demands and deliveries in $\{1,K\}$, including arbitrary tie choices.
The main integer guarantee and sharpness family have been checked in Lean
4.34.0 and replayed by the independent Nanoda checker. The searches and
implementation checks are recorded below.
Literature priority remains unconfirmed.

\medskip\noindent\textbf{Result.} If $L_{\rm init}$ is the initial assignment-LP
value, the algorithm returns capacity at most $L_{\rm init}+K-2$.
Since $L_{\rm init}\le\OPT$ and any present large delivery forces $\OPT\ge K$,
this proves the conjectured factor two. More precisely, for each integer
$K\ge2$ the exact worst-case approximation ratio, with original-index tie
breaking, is $2-2/K$. The additive bound is also attained for every such $K$.
Thus the best constant valid uniformly in $K$ is two, approached as
$K\to\infty$. If every delivery has size one, all delivery orders
coincide and are optimal.

\section{The target}
Let $K>1$ be an integer, $x_i\in\{1,K\}$, $y_i$ positive integers, and
$\sum_i x_i=\sum_i y_i$. Demand order is fixed. For a delivery order $q$, put
\[
s_0=0,\qquad s_j=\sum_{i=1}^j(q_i-y_i),\qquad
C(q)=\max_{1\le j\le n}(s_j+y_j)-\min_{0\le j\le n}s_j.
\]
Thus both delivery peaks and consumption troughs count. The initial inventory
can be shifted by $-\min_j s_j$; $s_n=0$ makes this a cyclic schedule.

The iterative rounding algorithm fixes positions in order. At each position it
tries each remaining delivery, solves the assignment LP with this additional
fixing, and keeps a candidate with the smallest LP value. Ties follow original
input index. Conjecture 3.1.1 on printed page 21 of
\href{https://perso.limos.fr/~lulorieau/docs/thesis/Master_Thesis.pdf}{Lorieau's thesis}
asserts $C(q_{\rm IR})\le2\OPT$. The unrestricted one-dimensional
conjecture remains outside this result.
The factor-two theorem for a different greedy algorithm in that thesis does
not by itself imply this conjecture.

\section{Eliminating the assignment matrix}
Write $h=K-1$. After an integral prefix has been fixed, suppose $r$ slots and
$m$ large deliveries remain.
There is also a direct normalization: replace each delivery $q_j$ by
$q_j-1\in\{0,h\}$ and each demand by $d_j=y_j-1\ge0$. Post-consumption
inventories remain unchanged, while every pre-consumption peak decreases
by exactly one. Hence the original capacity is exactly one plus the capacity
in this zero-or-$h$ model. This shift preserves all comparisons made
by the LP-based algorithm.
\begin{lemma}
The remaining fractional assignment is equivalent to choosing
$q_j=1+h p_j$, where $0\le p_j\le1$ and $\sum_jp_j=m$.
\end{lemma}
\begin{proof}
Given a doubly stochastic assignment matrix, let $p_j$ be the sum of entries
in column $j$ over rows for large deliveries. Column sums and row sums give
the stated bounds and total. Conversely, when $0<m<r$, put $p_j/m$ in each
large row of column $j$, and $(1-p_j)/(r-m)$ in each small row. Every row and
column sums to one and the column's delivery is $1+h p_j$. For $m=0$ or $m=r$,
use the uniform matrix on the remaining identical deliveries. Fixed prefix
rows and columns are kept integral.
\end{proof}
Consequently the LP asks for a minimum inventory range with $q_j$ fixed in
the prefix and $q_j\in[1,K]$ afterwards, subject to $\sum_jq_j=\sum_jy_j$.

\section{Equivalence with the assignment-LP rule}
The file \texttt{AssignmentLP.lean} defines the residual LP independently
of the closed-form score. Its predicate \texttt{MatrixBand} consists of a
nonnegative doubly stochastic assignment matrix, the inventory peak and
trough inequalities, and the bounds imposed by the already fixed prefix.
\texttt{LPValue} asserts both attainment and a lower bound for every feasible
assignment. The theorem \texttt{matrixLP\_value} identifies this minimum with
\texttt{prefixValue}+1; the added one restores original units.

In \texttt{Algorithm1.lean}, \texttt{Algorithm1Run} compares these independently
defined LP optima. The theorem \texttt{run\_iff\_algorithm1} proves that this
rule and the score-based \texttt{Run} have exactly the same complete
runs on every valid prefix. Validity requires only nonnegative normalized
demands, correct remaining length and supply total, and consistent prefix
extrema; it does not assume an LP formula. Candidate feasibility and minimum
attainment are consequences of the preceding theorems.

The corollaries \texttt{algorithm1\_additive\_guarantee} and
\texttt{algorithm1\_ratio\_bound} directly transfer the bounds to the LP rule.
\texttt{algorithm1\_two\_approximation} includes the all-small case, while
\texttt{algorithm1\_exists} proves existence of a run and an actual integral
schedule of its stated capacity. The theorem \texttt{sharp\_lp\_run} in
\texttt{Result.lean} transfers the entire extremal family as well.
All minimum-LP ties are allowed; the first-index rule is one permitted rule.
This does not assert that all tie paths occur for a single fixed input order.

\paragraph{Source convention.}
Figure 1.1 on printed page 4 of Lorieau's thesis displays row and column
bounds with $\le1$, whereas the text before Proposition 3.1.2 on printed
page 16 explicitly calls the feasible matrices doubly stochastic.
We formalize the standard assignment interpretation with row and column sums
\emph{equal} to one and all supplies used. We do not assert equivalence with
the literal weaker inequalities printed in Figure 1.1. Deleting fixed
permutation rows and columns yields our residual matrix; adjoining them
recovers the complete assignment. The relationship to the source's notation,
including this explicit convention, remains open to mathematical review.

\section{An exact certificate for the interval relaxation}
More generally let $q_j\in[\ell_j,u_j]$, with
$\sum_j\ell_j\le\sum_j y_j\le\sum_j u_j$. Index post-consumption inventory
vertices cyclically by $0,\ldots,n-1$, identifying $s_n=s_0$.
For each day $j$ insert two directed edges:
\[
(j-1)\longrightarrow j\quad\text{of weight }\ell_j-y_j,
\qquad
j\longrightarrow(j-1)\quad\text{of weight }y_j-u_j.
\]
An edge $a\to b$ of weight $w$ means $s_b\ge s_a+w$.
Let $d(i,j)$ be the maximum weight of a directed path from $i$ to $j$,
including the empty path when $i=j$, and let $p_j=\max_i d(i,j)$.

\begin{proposition}
The interval relaxation has exact value
\[
L=\max_{1\le j\le n}(p_j+y_j).
\]
A feasible solution is $q_j=p_j-p_{j-1}+y_j$, with cyclic indices.
The formula uses only integer arithmetic when the inputs are integral.
\end{proposition}
\begin{proof}
The assumptions ensure a feasible choice of supplies with the required total,
so all directed cycles have nonpositive weight. Maximum path weights are
finite and may be taken over simple paths. Every edge implies
$p_b\ge p_a+w$, so the displayed $q_j$ satisfies all supply bounds.
Its total telescopes to $\sum_jy_j$. Since empty paths give $p_j\ge0$,
the inventory potential $p$ has troughs at least zero and peaks at most $L$;
its capacity is at most $L$.

For the reverse inequality, every feasible inventory satisfies
$s_j-s_i\ge d(i,j)$. Capacity is at least the difference between the peak
of day $j$ and any trough $i$, hence at least $y_j+d(i,j)$. Maximizing gives
$C\ge L$. Translating the feasible potential to make $s_0=0$ preserves capacity.
\end{proof}
On a cycle, simple paths are clockwise or counterclockwise arcs. Thus this
certificate can be computed in $O(n^2)$ operations. It certifies both a lower
bound (one maximizing path) and a matching feasible fractional assignment.
No numerical LP tolerance is part of an algorithmic decision.

\section{Optimality for congruent demands}
\begin{proposition}
If every demand satisfies $y_j\equiv1\pmod{K-1}$, iterative rounding is
optimal, for every input order and every tie choice. In particular, it is
optimal for $K=2$ and arbitrary positive integer demands.
\end{proposition}
\begin{proof}
After any admissible fixed prefix, each lower and upper supply bound is
either $1$ or $K$. All edge weights in the preceding graph are consequently
multiples of $h=K-1$. The constructed potentials $p_j$ are multiples of $h$,
and $q_j=p_j-p_{j-1}+y_j$ is congruent to $1$ modulo $h$.
For an unfixed slot it lies in $[1,K]$, so it must be either $1$ or $K$.
The total-supply equality supplies exactly the remaining number of large
deliveries. Thus every prefix LP has an integral optimal completion.

At each step the LP score of a candidate therefore equals the cost of its
best integral completion. Taking a minimum over the remaining candidates
preserves the best completion cost of the current prefix. Induction from
the empty prefix proves that the returned order is optimal. For $K=2$, the
congruence assumption modulo $1$ is automatic.
\end{proof}
For $K\ge3$, integral demands need not satisfy this congruence.
The following sections treat those cases.

\section{A completion bound}
\begin{lemma}
For any feasible fixed prefix, the optimum integral completion has capacity
at most $L+(K-1)$, where $L$ is the fractional optimum with that prefix.
\end{lemma}
\begin{proof}
For a fractional solution write $P_j=\sum_{i\le j}p_i$ over the remaining slots.
Replace $p_j$ by $b_j=\lfloor P_j\rfloor-\lfloor P_{j-1}\rfloor$.
Because $0\le p_j\le1$, we have $b_j\in\{0,1\}$; because $P_r=m$ is integral,
$\sum_jb_j=m$. The cumulative delivery error is
$h(\lfloor P_j\rfloor-P_j)\in(-h,0]$. This shifts both the peak and trough
at day $j$ by the same error. The range increases by at most $h$ and the
fixed prefix is unchanged.
\end{proof}
For integral data the bound sharpens to $L+K-2$: use the integer supplies
from the path certificate. The cumulative rounding error is then one of
$0,-1,\ldots,-(K-2)$, rather than an arbitrary real number in $(-h,0]$.
This lemma alone does not control the iterative algorithm: its decisions need
not follow this particular completion. The next sections control every decision.

\section{An exact formula after a prefix}
We now work in the normalized model with deliveries $0$ or $h=K-1$ and
demands $d_j=y_j-1\ge0$. Let a fixed prefix have current inventory $s$,
minimum inventory $A$, and maximum inventory $B$, including both types of
event and the initial zero. Thus $A\le0\le B$ and $A\le s\le B$.
Suppose the remaining demands are $e_1,\ldots,e_r$. The total required
remaining supply $\sum e_i-s$ lies in $[0,rh]$.
Define, with empty maxima interpreted as zero,
\begin{align*}
D&=\max\left(0,\max_{1\le j\le r}\left[\sum_{i=1}^j e_i-hj\right]\right),\\
T&=\max\left(0,\max_{1\le i\le r}\left[\sum_{j=i}^r e_j-h(r-i)\right]\right),\\
V&=\max\left(0,\max_{1\le i\le j\le r}
             \left[\sum_{t=i}^j e_t-h(j-i)\right]\right).
\end{align*}
\begin{lemma}\label{lem:prefix}
The optimum fractional completion has capacity
\[
\max\left(V,\ \max(B,T)-\min(A,s-D)\right).
\]
\end{lemma}
\begin{proof}
Consider a proposed inventory band $[\alpha,\beta]$ containing the prefix,
so $\alpha\le A$ and $\beta\ge B$. Write $E_j=\sum_{i=1}^j e_i$.
The reachable interval after the next $j$ demands has endpoints obeying
\[
a_0=b_0=s,\quad
a_j=\max(\alpha,a_{j-1}-e_j),\quad
b_j=\min(\beta-e_j,b_{j-1}+h-e_j),
\]
provided no interval has become empty. This recurrence includes the
pre-consumption peak constraint, not just the post-consumption inventory.
Using $e_j\ge0$ gives
\begin{align*}
a_j&=\max(\alpha,s-E_j),\\
b_j&=\min\left(s+hj-E_j,
       \min_{1\le i\le j}\left[\beta-(E_j-E_{i-1})+h(j-i)\right]\right).
\end{align*}
These displayed intervals are nonempty for every $j$ precisely when the
following comparisons hold: $\beta-\alpha\ge V$, $\alpha\le s-D$, and
$s\le\beta$. To see sufficiency as well as necessity, compare both terms
in the maximum for $a_j$ with every term in the minimum for $b_j$;
the remaining comparisons follow from $e_i\ge0$ and $h\ge0$.
Their nonemptiness then justifies the reachable-interval recurrence by induction.

The terminal condition $0\in[a_r,b_r]$ additionally requires
$\alpha\le0$, $\beta\ge T$, and $0\le E_r-s\le rh$.
The prefix already gives $\alpha\le A\le0$ and $s\le B\le\beta$;
the total-supply condition gives the last requirement. Thus feasibility is
equivalent to
\[
\beta-\alpha\ge V,\qquad
\alpha\le\min(A,s-D),\qquad \beta\ge\max(B,T).
\]
Minimizing $\beta-\alpha$ under these three conditions proves the formula.
For $r=0$, feasibility requires $s=0$ and the formula reduces to $B-A$.
\end{proof}

\section{The two choices cannot both cross the bound}
Let $L_0=L_{\rm init}-1$ be the initial LP value in the normalized model.
Suppose both delivery sizes remain available. Let $d$ be the next demand,
and let $D,T,V$ be computed from demands \emph{strictly after} that day.
For a tentative next fractional delivery $z\in[0,h]$, define
\[
R=d+D,\quad B_* =\max(B,T),\quad
F=\max(V,B_*-A,R),\quad
u=B_*+R-s,\quad v=s-A.
\]
There is at least one large and one small delivery left; consequently every
$z\in[0,h]$ leaves a feasible total fractional supply for the later slots.
After that delivery the prefix extrema and state are
$\min(A,s+z-d)$, $\max(B,s+z)$, and $s+z-d$.
Lemma~\ref{lem:prefix} therefore gives its exact score as
\begin{equation}\label{eq:score}
f(z)=\max(F,u-z,v+z).
\end{equation}
Here $v\le F$, since $s\le B\le B_*$.
If $L$ is the current LP value before fixing the next slot, then
$L=\min_{0\le z\le h} f(z)$, and hence
\begin{equation}\label{eq:lower}
F\le L,\qquad u-h\le L.
\end{equation}

We also have $R\le L_0$. Indeed, $R$ is the maximum, over all intervals
beginning at the next day, of
\[
\sum_{i=a}^b d_i-h(b-a).
\]
For any feasible initial fractional schedule, its peak at $a$ minus its
trough at $b$ equals
$\sum_{i=a}^b d_i-\sum_{i=a+1}^b z_i$, which is at least the displayed
quantity. Thus each such quantity is a lower bound on the initial LP value.

\begin{theorem}\label{thm:main}
For positive integer demands and $x_i\in\{1,K\}$ with integer $K\ge2$,
every run of iterative rounding satisfies
\[
C(q_{\rm IR})\le L_{\rm init}+K-2\le\OPT+K-2.
\]
In particular it is a 2-approximation. All choices among exactly tied
minimum LP scores are permitted.
\end{theorem}
\begin{proof}
Set $H=L_0+h-1$, an integer. We prove by induction that the normalized
LP value after every fixing is at most $H$. The empty prefix satisfies
this because $h\ge1$. If only one delivery size remains, the total-supply
constraint already forces every remaining fractional delivery to that
size, so the LP value cannot change.

Otherwise use \eqref{eq:score} and suppose the current value $L\le H$.
If both endpoint scores exceeded $H$, then
\[
f(0)=\max(F,u)>H\quad\Longrightarrow\quad u\ge H+1,
\]
while $F\le H$ and $u-h\le H$ from \eqref{eq:lower} give
\[
f(h)=\max(F,u-h,v+h)>H\quad\Longrightarrow\quad v+h\ge H+1.
\]
Integrality is used in these two implications. Adding yields
$u+v+h\ge2H+2$. But
\[
u+v+h=(B_*-A)+R+h
\le F+L_0+h\le H+L_0+h=2H+1,
\]
a contradiction. At least one endpoint score is at most $H$, so every
minimum-score choice preserves the induction bound.

After the final fixing, the LP value is the actual normalized capacity.
Adding back the baseline one gives
$C(q_{\rm IR})\le H+1=L_{\rm init}+K-2$.
The initial LP is a relaxation, so $L_{\rm init}\le\OPT$.
If a large delivery is present, $K\le\OPT$, and hence
$C(q_{\rm IR})\le2\OPT-2<2\OPT$.
If none is present, all deliveries are identical and the algorithm is optimal.
\end{proof}

\noindent\textbf{Real-valued extension.}
The same argument without the integer improvement uses $H=L_0+h$:
both endpoint scores above $H$ would imply $u+v+h>2H$, whereas the upper
bound is $2H$. Thus, for real $K>1$ and real $y_i\ge1$, the bound is
$C(q_{\rm IR})\le L_{\rm init}+K-1$. Demands below one are outside this
argument: their normalization would be negative, invalidating the reachability
formula used above.

\section{Sharpness for every integer delivery size}
\begin{theorem}\label{thm:tight}
For each integer $K\ge3$ there is an input with $n=2K-3$ such that
original-index iterative rounding has
\[
L_{\rm init}=\OPT=K,\qquad C(q_{\rm IR})=2K-2.
\]
Consequently both the additive bound in Theorem~\ref{thm:main} and the
approximation ratio $2-2/K$ are sharp. For $K=2$ the algorithm is always
optimal, so the same conclusion about the ratio holds.
\end{theorem}
\begin{proof}
Put $h=K-1\ge2$. In the normalized zero-or-$h$ model, take the input
delivery order
\[
z=(h,0)^{h-1},h,
\]
where a superscript means repetition of the pair. Define the demand sequence by
\[
d_1=1,\qquad
(d_{2j},d_{2j+1})=(h-1-j,j+1)\quad(1\le j\le h-2),
\qquad d_{2h-2}=h,\quad d_{2h-1}=h-1.
\]
The paired portion is empty when $h=2$. There are $h$ large deliveries and
$h-1$ zero deliveries. The total demand is
$1+(h-2)h+h+(h-1)=h^2$, equal to total supply.
Adding one to all entries gives valid positive integer input lists $X,Y$.

The alternative delivery order
\[
z^*=(0,h)^{h-1},h
\]
has every inventory level in $[-1,h-1]$: the initial demand leaves state
$-1$; each intermediate pair of demands sums to $h$; the penultimate demand
is $h$; and the final day returns the state to zero. Its normalized capacity
is $h$, and its original capacity is $K$. The demand $d_{2h-2}=h$ is itself
a lower bound of $h$ on any normalized LP capacity. Hence
$L_{\rm init}=\OPT=K$ in the original model.

We now verify that the algorithm preserves the displayed input order.
All demands lie in $[0,h]$, so every suffix has $D=0$. For every nonempty
suffix, $T=h-1$: its last demand is $h-1$, and every preceding demand is
at most $h$. Until the penultimate day is consumed, a suffix containing
that day has $V=h$. Substituting these quantities into
\eqref{eq:score} gives the following exact normalized LP scores.
\[
\begin{array}{c|c|c|c}
\text{day}&f(0)&f(h)&\text{chosen delivery}\\\hline
1&h&h&h\\
2j\ (1\le j\le h-2)&h+j-1&2h-1&0\\
2j+1\ (1\le j\le h-2)&h+j&h+j&h\\
2h-2&2h-1&2h-1&0\\
2h-1&\text{unavailable}&2h-1&h
\end{array}
\]
For completeness, the prefix before day $2j$ has
\[
s=h-1,\quad A=0,\quad B=h+j-1.
\]
Its zero delivery and demand $h-1-j$ leave $s=j$. The next large delivery
raises the peak to $h+j$ and its demand $j+1$ restores $s=h-1$,
proving these states inductively. Before day $2h-2$, the state is $h-1$
and the extrema are $A=0$, $B=2h-2$. Its zero delivery leaves state $-1$;
the final large delivery returns to zero without exceeding the preceding peak.
Every tie in the table is resolved as displayed because the next member of
the input sequence has the smallest remaining original index.

Thus the returned normalized trace has minimum $-1$ and maximum $2h-2$.
Its normalized capacity is $2h-1$, giving original capacity
$2h=2K-2$. Dividing by $\OPT=K$ proves the claimed lower ratio;
Theorem~\ref{thm:main} supplies the matching upper ratio. When $K=2$,
that theorem gives $C(q_{\rm IR})\le\OPT$ directly.
\end{proof}

\noindent\textbf{A small example.} For $K=5$, the construction gives
\[
\begin{aligned}
X=q_{\rm IR}&=(5,1,5,1,5,1,5),\\
Y&=(2,3,3,2,4,5,4),\\
q^*&=(1,5,1,5,1,5,5).
\end{aligned}
\]
The returned capacity is $8$, while the optimum is $5$.

\noindent\textbf{What is sharp.}
For fixed $K$, the family attains the ratio $2-2/K$.
Across all integer $K$, the supremum is two: no smaller uniform constant
works, although no finite instance with a large delivery has ratio exactly two.
The construction uses exact LP ties and the specified original-index rule.
It does not establish the same lower bound for every alternative tie-breaking
rule, and it is not an impossibility result for other algorithms.

\section{Computing the same decisions without LP solves}
The suffix quantities admit backwards recurrences. Index days from $1$ to
$n$ and initialize $D_{n+1}=T_{n+1}=V_{n+1}=0$. Then
\begin{align*}
D_t&=\max(0,d_t-h+D_{t+1}),\\
T_t&=\max\left(T_{t+1},\sum_{j=t}^n d_j-h(n-t)\right),\\
V_t&=\max(V_{t+1},d_t+D_{t+1}).
\end{align*}
The initial normalized LP value is $\max(V_1,T_1+D_1)$.
At day $t$, use $R=d_t+D_{t+1}$, $T=T_{t+1}$, and $V=V_{t+1}$ in
\eqref{eq:score}. Maintain $A,B,s$ and two queues of original indices,
one for each delivery size. Each day compares at most two exact scores
and pops the appropriate index. This takes $O(n)$ arithmetic operations
and $O(n)$ storage, including preprocessing, while preserving the specified
original-index tie rule. Integer bit complexity also depends on input magnitude.
The implementation is in \texttt{fast\_algorithm.py}; its fixed-example
cross-checks against the cycle-path oracle are in \texttt{verify\_examples.py}.

\section{Finite experiment and tie handling}
The search independently enumerates binary delivery orders to obtain $\OPT$,
and follows every candidate tied for the best exact LP value to obtain the
largest possible iterative-rounding cost. Every such path is realizable by
the stated input-index tie rule: take its resulting delivery sequence as
the input order. At each step the next chosen delivery then has the first
remaining original index.

The initial experiment also probed two stronger statements as possible proof guides:
$C(q_{\rm IR})\le L_0+K-1$, and whether all unshifted inventory levels lie in
$[-\OPT,\OPT]$. The origin-centered band statement is false. For example,
$K=3$, $Y=(2,3,2)$ and input $X=(3,3,1)$ permit the returned order $(3,3,1)$,
whose peak is $4$, although $(1,3,3)$ has capacity $3=\OPT$.
The additive bound was not falsified and is strengthened by Theorem~\ref{thm:main}.

The Azure run checked 7,094,426 input trials and 47,073,459 tied algorithm
paths in 240 seconds. The first 122,861 trials exhaust all positive-demand
inputs with $2\le n\le6$, $2\le K\le5$ and both supply types present;
the rest are random trials with $7\le n\le12$ and $2\le K\le25$ and may
include repeats. No factor-two counterexample was found. The largest observed
ratio was $12/7$, with $K=7$, $\OPT=7$, and algorithm cost $12$.

A separate full assignment LP implementation agreed with 300 exact oracle
evaluations. After deriving \eqref{eq:score}, 108 candidate scores on the nine
record-setting examples were checked against the independent cycle-path oracle.
Heavy enumeration ran only in Azure with an 8 GiB memory cap, a 12-minute
worker limit, and an independent deletion guard. Both resource groups and all
job resources were confirmed deleted at 13:02:59 UTC on 1 October 2026.
The sharpness construction was derived algebraically, without an additional
search. Five fixed instances at $K\in\{3,4,5,7,16\}$ were checked locally
against the independent cycle-path oracle, including all 105 candidate scores
and the claimed comparison schedules.

\section{Formal verification and reproducibility}
The original nine source modules were compiled in a clean Azure environment using Lean
4.34.0 and mathlib commit
\texttt{5ed2965256430c3649e86755f9576b54eca72435}.
The checked development is structured as follows:
\begin{itemize}
\item \texttt{Assignment.lean}: both directions of the fractional assignment
reduction, including empty delivery classes.
\item \texttt{Reachability.lean}: actual real-valued fractional deliveries,
exact band feasibility, and attainment of the prefix LP value.
\item \texttt{IntegerModel.lean}: casting the real model to exact integer
scores, the score formula, and the bound for the two choices.
\item \texttt{Rounding.lean}: all permitted minimum-score runs, their
existence, and the induction proving the additive guarantee.
\item \texttt{Normalization.lean}: equivalence of original and normalized
band feasibility, with a one-unit shift of the peak bound.
\item \texttt{Schedules.lean}: arbitrary integral comparison schedules and
the ratio inequality $(h+1)(C+1)\le2h(b-a+1)$.
\item \texttt{Ties.lean}: realization of each permitted tie path by an
original input order whose next item is a minimum-score candidate.
\item \texttt{Sharpness.lean} and \texttt{Result.lean}: the parameterized
family, its bad run, an optimal integral schedule, and the matching lower bound.
\end{itemize}
Here $h=K-1$, $C$ is normalized algorithm capacity, and $b-a$ is the normalized
capacity of an arbitrary comparison schedule. The formal sharpness parameter
$r\in\mathbb N$ represents $K=r+3$; it is not a finite enumeration.

The original twelve audited target statements and their complete dependencies were
exported and checked by unmodified Nanoda commit
\texttt{3a2407216ee84a75f9e1aead6803d0578be06ae7}.
Nanoda accepted 9451 declarations. The only axioms in the export are the
standard \texttt{propext}, \texttt{Classical.choice}, and \texttt{Quot.sound}.
There are no custom axioms, admitted proofs, unsafe or partial declarations
in the exported dependency set. No native-decision proof shortcuts are used.
Lean rejected an invalid proof of \texttt{False}; Nanoda accepted a valid
control and rejected the same exported control with its claimed type changed
to \texttt{False}. No kernel or checker source was modified. The exporter's
only working-tree change selected the pinned Lean toolchain.

These checks establish the formal statements and their dependencies; they
are not an external peer review of the problem translation or a priority
certification. The Python and JavaScript implementations were checked
separately. The auxiliary
cycle-path certificate, completion-rounding observation, real-valued extension,
and arithmetic complexity analysis are not claimed as separately formalized
theorems in this package.

The clean audit ran in Azure job \texttt{01bb7023795d4c02}, with a 12 GiB
memory limit, a 35-minute worker limit and an independent 45-minute compute
deletion guard. Source hashes, full compiler diagnostics, checker configuration,
a compressed proof export and negative-control logs accompany the package.

On 4 October 2026, Azure job \texttt{c81bbb49c1b84185} compiled all eleven
current proof modules, including the two new bridge modules, and checked the
axiom dependencies of twenty-four target statements. They use only
\texttt{propext}, \texttt{Classical.choice}, and \texttt{Quot.sound}.
The new bridge is not included in the historical Nanoda export: its present
verification is by Lean. Full diagnostics, exact source hashes, a rejected
negative control, and resource-cleanup evidence accompany the new report
\texttt{evidence/algorithm-bridge-audit.json}.

\section{An initial experiment with alternative tie rules}
The tight family can behave better under another tie rule. However, always
preferring the small delivery on an exact tie does not make the algorithm
optimal. A bounded Azure experiment with seed 20261001 checked 10000 random
inputs with $3\le n\le10$ and $3\le K\le16$, allowing repeats. Each optimum
was obtained by enumerating all binary orders; every new ratio record was
checked against the independent cycle-path LP oracle.

For the small-first rule, one recorded input is
\[
X=(1,1,1,5,5,5),\qquad Y=(2,2,3,5,4,2).
\]
The returned order $(1,1,5,1,5,5)$ has capacity $7$, whereas
$(5,1,1,5,5,1)$ has capacity $5$, which is optimal because a delivery of size
$5$ is present. This refutes exact optimality for that rule. It does not
settle whether the rule has a better worst-case approximation factor.
The interactive laboratory includes both this example and the sharpness family.

\section{Attribution}
Prepared by MathIsEvenEasier with OpenAI Codex (GPT-6 Astra).
Research inspired by \texttt{@xamualexander}, Dr. Samuel Allen Alexander.
Still there.


\end{document}
